\scanpage{33}
% @ 1
E\uline{xemples de groupes algébriques affines sur} $k$.

% @ 2–8
\diagram{\scriptsize\(\begin{array}{@{}l@{\;}l@{\;}l@{\;}c@{\;}l@{}}
\mathbb{G}_m, & \mathbb{G}_a, & \mathrm{GL}(n), & \mathrm{O}(Q)., &
 \text{ss-groupe de }\mathrm{GL}(n)\ \text{invariant des tenseurs}\\[-.5ex]
&&\mathrm{SL}(n)&\begin{array}{c}\cap\\[-.8ex]\mathrm{GL}(n)\end{array}
 &\hfill .\,.\,.\,.\,. \\
&&\mathrm{Tr}(n)&&\\
\mathrm{Sp}.&&\mathrm{Tr}_{\unsure{0|o}}(n)&&\\
\multicolumn{2}{r}{\mathrm{Diag}_0(n)} & \multicolumn{1}{l}{\subset\mathrm{Diag}(n)} &
\multicolumn{2}{l}{\text{immersion de }\mathbb{G}_a\text{ dans }\mathbb{G}\mathrm{L}(2)}\\
\multicolumn{2}{c}{\vcenter{\hbox{\tikz[baseline]\node[rotate=90,inner sep=0pt]{\(\simeq\)};}}} &&&\\
\multicolumn{2}{c}{\mathbb{G}_m}&&&
\end{array}\)}

% @ 9
Soit $M$ un $k$-module.
% @ 10–11
\[
 W(M):\unsure{\{|\varepsilon}\ k'\longmapsto M\otimes_k \unsure{k|K}'
 \qquad\text{\parbox{.4\linewidth}{(en fait un $\mathbb{O}_k$-Module, pas seulement un groupe)}}
\]

% @ 11–13
$\cdot$ Pour que ce soit représentable, il faut que $W(M)$ transforme injections $k'\mathrel{\tikz[baseline=-.3ex,x=1em,y=1em]{\draw[line width=.45pt] (0,0)--(1,0)--(1,.12);}} k''$ en injections, donc que $M$ soit p\uline{lat} sur $k$. \deleted{\unsure{So|Si}}

% @ 14–15
Il faut qu'il transforme produits (infinis) d'algèbres en produits
\[
 M\otimes_k\prod k'_i\longrightarrow\prod_i M\otimes_k k'_i
 \qquad\text{bijectif.}
\]

% @ 16–18
\deleted{On voit que ceci} faisant tous les $k'_i$ égaux : $k$, on trouve
\[
 M\otimes k^I\xrightarrow{\sim}\deleted{\unsure{\prod_i|\prod}} M^I
\]
% @ 19–21
Et on voit tout de suite que cela implique que $M$ est de type fini, donc \deleted{\unsure{un|une}} quotient d'un \tikz[baseline=(s.base)]{\node[inner sep=1pt](s){libre \deleted{\unsure{fini|f}}};\node[anchor=south,inner sep=1pt] at ([yshift=4pt]s.north){\deleted{$R\mathrel{\tikz[baseline=-.3ex,x=1em,y=1em]{\draw[line width=.45pt] (0,0)--(1,0)--(1,.12);}}L$}};\fill ([xshift=12pt,yshift=-2pt]s.south) circle (.4pt);} \deleted{\unsure{librement|comme}} $0\mathrel{\tikz[baseline=-.3ex,x=1em,y=1em]{\draw[line width=.45pt] (0,0)--(1,0)--(1,.12);}}R\mathrel{\tikz[baseline=-.3ex,x=1em,y=1em]{\draw[line width=.45pt] (0,0)--(1,0)--(1,.12);}}L\mathrel{\tikz[baseline=-.3ex,x=1em,y=1em]{\draw[line width=.45pt] (0,0)--(1,0)--(1,.12);}}M\mathrel{\tikz[baseline=-.3ex,x=1em,y=1em]{\draw[line width=.45pt] (0,0)--(1,0) (1,-.07)--(1,.09);}}0$
% @ 22–25
et comme $M$ est plat, cel reste \unsure{ex.ct|exact} par $\otimes_k k^I$, donc $R\otimes k^I\overset{\sim}{\mathrel{\tikz[baseline=-.3ex,x=1em,y=1em]{\draw[line width=.45pt] (0,0)--(1,0)--(1,-.07) (1,0) to[out=130,in=-30] (.85,.23) to[out=150,in=130] (.94,.26) to[out=-45,in=90] (1,0);}}}R^I$, donc $R$ de type fini, donc $M$ de présentation finie, donc $M$ plat de prés. finie, donc $M$ projectif de t-f. (facteur direct d'un libre)

% @ 26–27
Inversement, si $M$ projectif de t.f. i.e.. facteur direct d'un libre \tikz[baseline=(s.base)]{\node[inner sep=1pt](s){$L$};\node[anchor=south west,inner sep=1pt](a) at ([xshift=-2pt,yshift=3pt]s.north east){\inserted{$\simeq k^{\unsure{n}}$}};\draw ([xshift=2pt,yshift=-5pt]s.south east) to[out=90,in=210] ([yshift=-2pt]a.south east) to[out=35,in=-90] ([xshift=2pt,yshift=2pt]a.south east);}, par projecteur $p$, on a
% @ 28–29
\[
 \begin{aligned}
 M\otimes_k k'&=\mathrm{Ker}\big((1-p)_{k'}:k'^{n}\to k'^{n}\big)\\
 W(M)&=\mathrm{Ker}\big(\mathbb{E}^{n}\overset{1-p}{\mathrel{\tikz[baseline=-.3ex,x=1em,y=1em]{\draw[line width=.45pt] (0,0)--(1,0)--(1,.12);}}}\mathbb{E}^{n}\big)
 \end{aligned}
\]
% @ 30
c'est repr.

% archive sheet number: 32
